%------------------------------------------------------------------------------%
\section{Funções}

%------------------------------------------------------------------------------%
\begin{frame}{Funções}
  Podemos agora explorar diversas funções e seus efeitos no jogo
\end{frame}

%------------------------------------------------------------------------------%
\begin{frame}{Operações Suportadas}
  \renewcommand{\arraystretch}{1.8}
  \centering
  \begin{tabular}{c<{\quad}l<{\quad}>{\quad}c<{\qquad}}
    \code{+}                & soma      &  \code{3 + 4}  \\ \pause
    \code{-}                & subtração &  \code{3 - 4}  \\ \pause
    \code{*}                & produto   &  \code{3 * 4}  \\ \pause
    \code{/}                & divisão   &  \code{3 / 4}  \\ \pause
    \code{\textasciicircum} ou \code{**} & potência  & \code{3 ** 4}
  \end{tabular}
\end{frame}

%------------------------------------------------------------------------------%
\framedgraphic{Nova Velocidade}{vel-power-2}{}
\framedgraphic{Nova Velocidade}{vel-power-4}{}

%------------------------------------------------------------------------------%
\begin{frame}{Algumas Funções Suportadas}
  \footnotesize
  \begin{columns}[t]
    \column{0.3\linewidth}
    Raiz quadrada e módulo

    \medskip
    \code{raiz(a)} ou \code{sqrt(a)}

    \medskip
    \code{módulo(a)} ou \code{abs(a)}

    \medskip
    Mínimo e máximo

    \medskip
    \code{min(a,b)}

    \code{max(a,b)}

    \medskip
    Exponenciais e Logaritmos

    \medskip
    \code{exp(a)}

    \code{ln(a)} ou \code{log(a)}

    \code{log10(a)}

    \column{0.33\linewidth}
    Trigonométricas

    \bigskip
    \code{sen(a)} ou \code{sin(a)}

    \code{cos(a)}

    \code{tg(a)} ou \code{tan(a)}

    \bigskip
    \code{arcsen(a)} ou \code{arcsin(a)}

    \code{arccos(a)}

    \code{arctg(a)} ou \code{arctan(a)}

    \column{0.36\linewidth}
    Hiperbólicas

    \bigskip
    \code{senh(a)} ou \code{sinh(a)}

    \code{cosh(a)}

    \code{tgh(a)} ou \code{tanh(a)}

    \bigskip
    \code{arcsenh(a)} ou \code{arcsinh(a)}

    \code{arccosh(a)}

    \code{arctgh(a)} ou \code{arctanh(a)}

  \end{columns}
\end{frame}

%------------------------------------------------------------------------------%
\framedgraphic{Nova Velocidade}{vel-raiz}{}
\framedgraphic{Nova Velocidade}{vel-seno}{}

%------------------------------------------------------------------------------%
